\chapter{Myringotomy and Tube Insertion}

\section*{What Is This Surgery?}
Myringotomy with tympanostomy tube placement is a highly prevalent surgical procedure involving the creation of a small, precise incision in the tympanic membrane (eardrum) to facilitate the aspiration of accumulated middle ear fluid and the subsequent insertion of a tiny, hollow ventilation tube. This tube, commonly known as a grommet or pressure equalization (PE) tube, serves to maintain patency of the incision, thereby ventilating the middle ear space, equalizing pressure across the eardrum, and preventing the re-accumulation of fluid. The procedure is primarily indicated for the management of recurrent acute otitis media (AOM) and chronic otitis media with effusion (OME), conditions that can significantly impair hearing, impede speech and language development, and lead to other developmental challenges, particularly in the pediatric population. It represents a cornerstone intervention in pediatric otolaryngology, effectively restoring normal middle ear function and substantially improving patients' quality of life.

\section*{Alternative Names}
\begin{itemize}
  \item Ear tube surgery
  \item Tympanostomy tube placement
  \item Grommet insertion
  \item Ventilation tube insertion
  \item Pressure equalization (PE) tube insertion
  \item Myringotomy with PE tube
  \item Bilateral myringotomy and tubes (BMT)
\end{itemize}

\section*{Relevant Surgical Anatomy}
A thorough understanding of the external auditory canal (EAC) and middle ear anatomy is paramount for safe and effective myringotomy and tube insertion. The tympanic membrane (TM) forms the boundary between the EAC and the middle ear, comprising three layers: an outer epidermal layer, a middle fibrous layer (absent in the pars flaccida), and an inner mucosal layer. Key landmarks on the TM include the umbo (the most retracted point, where the malleus handle attaches), the malleus handle itself, and the cone of light (anteroinferiorly). The TM is divided into four quadrants (anterior-superior, anterior-inferior, posterior-superior, posterior-inferior) for descriptive purposes, with the anterior-inferior quadrant being the preferred site for incision due to its relative avascularity and distance from the ossicular chain and facial nerve.

The middle ear cavity is an air-filled space containing the three ossicles (malleus, incus, stapes) responsible for sound transmission. The Eustachian tube connects the anterior middle ear to the nasopharynx, regulating middle ear pressure and drainage. Important structures in close proximity include the facial nerve (coursing through the posterior-superior aspect of the middle ear), the chorda tympani nerve (a branch of the facial nerve, traversing the middle ear superior to the malleus and incus), and the jugular bulb (inferiorly). The arterial supply to the tympanic membrane and middle ear is primarily from branches of the maxillary artery (deep auricular artery) and the posterior auricular artery (stylomastoid artery). Venous drainage parallels the arterial supply, and lymphatic drainage typically follows the external jugular chain.

\section{section*{Surgical Principle \& Rationale}}
The fundamental surgical principle underlying myringotomy and tympanostomy tube insertion is to establish a temporary, artificial bypass for a dysfunctional Eustachian tube. The Eustachian tube's primary role is to ventilate the middle ear, equalize pressure with the ambient atmosphere, and drain secretions. When this tube fails to open and close properly, often due to inflammation, adenoid hypertrophy, or immaturity in children, a negative pressure environment develops within the middle ear. This negative pressure can lead to retraction of the tympanic membrane and the accumulation of fluid, resulting in otitis media with effusion (OME) or, if infected, acute otitis media (AOM).

The surgical intervention directly addresses these pathological processes by creating a direct, patent pathway for air exchange. The myringotomy incision allows for immediate aspiration of any existing middle ear fluid, thereby relieving pressure and removing the inflammatory milieu. The subsequent insertion of a tympanostomy tube maintains this opening, ensuring continuous ventilation of the middle ear space. This sustained aeration equalizes pressure, prevents the re-accumulation of fluid, and significantly reduces the frequency of acute infections by allowing the middle ear mucosa to return to a healthy, non-inflamed state. The tube essentially functions as a temporary, artificial Eustachian tube, providing a window for the natural Eustachian tube function to mature or recover over time.

\section{section*{History \& Surgical Evolution}}
The concept of incising the tympanic membrane to alleviate middle ear pathology dates back to the early 19th century. Sir Astley Cooper, a renowned British surgeon, is widely credited with performing one of the earliest documented myringotomies in 1801, primarily to relieve pressure and drain fluid from the middle ear. However, these initial procedures lacked a mechanism to maintain the patency of the incision, leading to rapid closure and frequent recurrence of symptoms. Various attempts were made over the subsequent decades to keep the incision open, including the use of small metal or rubber grommets, but none gained widespread acceptance or demonstrated consistent efficacy.

A pivotal breakthrough in the long-term management of middle ear ventilation came in 1954 when Dr. Beverly Armstrong, an American otolaryngologist, introduced the concept of inserting a small, hollow plastic tube into the myringotomy incision to maintain prolonged patency. His pioneering work with the "collar button" tube design revolutionized the treatment of chronic middle ear effusions and recurrent otitis media. Following Armstrong's innovation, numerous tube designs and materials have been developed, including various shapes (e.g., Shepard, T-tube, Paparella) and materials (e.g., silicone, fluoroplastic, titanium), each aiming to optimize retention time, minimize complications, and suit specific clinical scenarios. Today, myringotomy with tympanostomy tube insertion remains one of the most frequently performed surgical procedures in pediatric medicine globally, a testament to its enduring efficacy and safety.

\section{section*{Clinical Indications}}
Myringotomy and tympanostomy tube insertion are indicated for specific, well-defined criteria related to chronic or recurrent middle ear fluid and infection that have not responded to appropriate medical management.

\begin{itemize}
  \item Recurrent acute otitis media (AOM), typically defined as three or more distinct episodes within six months, or four or more episodes within 12 months, with at least one episode occurring in the preceding six months.
  \item Chronic otitis media with effusion (OME) lasting three months or longer, particularly when associated with documented bilateral hearing loss of 20 dB or greater, or unilateral hearing loss with significant impact.
  \item OME accompanied by significant speech or language delay, balance disorders, behavioral issues, or learning difficulties that are directly attributable to hearing impairment.
  \item Severe tympanic membrane retraction (atelectasis) or adhesive otitis media, which can predispose to the formation of cholesteatoma, a destructive skin cyst in the middle ear.
  \item Eustachian tube dysfunction causing persistent discomfort, recurrent barotrauma (e.g., in divers or frequent flyers), or chronic serous otitis media, especially in adults.
  \item Prevention of recurrent otitis media in children with craniofacial anomalies (e.g., cleft palate, Down syndrome) or other underlying conditions that significantly increase their risk for chronic middle ear disease.
  \item Complications of otitis media, such as impending mastoiditis, facial nerve paralysis, or intracranial complications, where rapid and sustained middle ear drainage and ventilation are critically important.
  \item As a diagnostic or therapeutic measure for specific middle ear conditions, such as perilymphatic fistula or Meniere's disease, though these are less common indications.
\end{itemize}

\section*{Clinical Positioning}
Myringotomy and tympanostomy tube insertion is generally considered a primary surgical treatment for specific, well-defined criteria of recurrent acute otitis media (AOM) and chronic otitis media with effusion (OME), particularly after a period of watchful waiting or a trial of medical therapy has proven ineffective or is deemed inappropriate. For recurrent AOM, it serves as a primary surgical intervention once a patient has met the established frequency criteria, often following multiple courses of antibiotics that have failed to prevent further infections. In the context of OME, it is a primary surgical option when the effusion persists for an extended duration (typically three months or more) and is associated with significant hearing loss, speech delay, or other developmental concerns, usually after an initial period of observation.

While antibiotics are the primary medical treatment for acute middle ear infections, tube insertion is positioned as a secondary intervention within the broader management strategy, aiming to prevent future infections and resolve chronic effusions by directly addressing the underlying ventilation deficit. In certain circumstances, such as severe tympanic membrane retraction, impending cholesteatoma, or specific craniofacial syndromes, tube placement may be considered a primary intervention to prevent more serious and irreversible complications. It is rarely a salvage procedure, but rather a proactive measure to mitigate the long-term sequelae of chronic middle ear disease.

\section*{Surgical Technique \& Procedure}
Myringotomy with tympanostomy tube insertion is typically performed as an outpatient procedure, usually lasting between 10 to 15 minutes per ear. In pediatric patients, general anesthesia is almost universally employed, while in cooperative adults, it can occasionally be performed under local anesthesia.

The general steps are as follows:
\begin{itemize}
  \item \textbf{Anesthesia and Patient Positioning:} For children, general anesthesia is induced, and the patient is positioned supine with the head turned to adequately expose the operative ear. The surgeon utilizes an operating microscope to provide magnified, stereoscopic visualization of the external auditory canal and tympanic membrane.
  \item \textbf{Ear Canal Preparation and Visualization:} The external ear canal is gently cleaned with an antiseptic solution, such as povidone-iodine or alcohol, taking care to avoid contact with the tympanic membrane. Any cerumen (earwax) or epithelial debris is meticulously removed using curettes or suction to ensure an unobstructed and clear view of the eardrum.
  \item \textbf{Myringotomy Incision:} Using a specialized, sharp myringotomy knife, the surgeon makes a small, radial incision, typically 2-3 mm in length, in the anterior-inferior quadrant of the tympanic membrane. This specific location is chosen to minimize the risk of injury to the ossicular chain and facial nerve, and it is an area that tends to heal well after tube extrusion.
  \item \textbf{Fluid Aspiration:} If middle ear fluid is present, it is carefully aspirated using a small, angled suction catheter. The fluid's consistency (serous, mucoid, or purulent) is noted, and a sample may be collected for microbiological culture if an active infection is suspected or if the patient is immunocompromised.
  \item \textbf{Tube Insertion:} A tympanostomy tube of appropriate size and design (e.g., a short-term Shepard grommet, a longer-term T-tube, or a collar button tube) is then carefully inserted through the myringotomy incision. The flanges of the tube are positioned to rest on either side of the tympanic membrane, securely holding it in place.
  \item \textbf{Post-Insertion Check and Antibiotic Instillation:} The surgeon confirms that the tube is properly seated, patent, and allows for free passage of air into the middle ear. The ear canal is then re-examined for any bleeding or residual debris. Topical antibiotic ear drops, often containing a fluoroquinolone or aminoglycoside, may be instilled into the ear canal at the conclusion of the procedure, particularly if purulent fluid was drained or as a prophylactic measure.
  \item \textbf{Closure:} No sutures are required for the myringotomy incision, as it is designed to self-seal around the tube. The patient is then awakened from anesthesia and transferred to the recovery area.
\end{itemize}

\section*{Preoperative Workup \& Preparation}
Preparation for myringotomy and tube insertion is designed to ensure patient safety, optimize surgical outcomes, and minimize anesthetic risks. A comprehensive preoperative workup typically involves several key steps.

\begin{itemize}
  \item \textbf{Medical History and Physical Examination:} A thorough review of the patient's medical history, including current medications, allergies, and previous anesthetic experiences, is conducted by both the referring physician and the otolaryngologist. A focused physical examination confirms the indications for surgery.
  \item \textbf{Audiometry and Tympanometry:} Preoperative audiometry is often performed, especially in cases of chronic OME, to objectively document the degree and type of hearing loss. Tympanometry is a crucial diagnostic tool that helps confirm the presence of middle ear effusion and assess Eustachian tube function.
  \item \textbf{Anesthesia Clearance:} For procedures requiring general anesthesia, a preoperative assessment by an anesthesiologist is mandatory. This evaluation ensures the patient is medically fit for anesthesia and may involve specific blood tests (e.g., complete blood count, coagulation profile), an electrocardiogram (ECG), or other cardiac evaluations, depending on the patient's age, comorbidities, and overall health status.
  \item \textbf{NPO Status:} Patients are strictly instructed to be NPO (nil per os, or nothing by mouth) for a specified period prior to general anesthesia, typically 6-8 hours for solid foods and 2-4 hours for clear liquids, to prevent the risk of pulmonary aspiration during induction.
  \item \textbf{Medication Adjustments:} Specific instructions are provided regarding which regular medications should be taken or held before surgery. Anticoagulants or antiplatelet agents may need to be temporarily discontinued or adjusted under medical supervision to minimize bleeding risk.
  \item \textbf{Informed Consent:} A detailed discussion with the patient or their legal guardians (parents) regarding the procedure, its anticipated benefits, potential risks, alternative treatments, and the expected postoperative course is conducted. Written informed consent is obtained prior to the procedure.
  \item \textbf{Prophylactic Antibiotics:} While not universally required, some surgeons may prescribe a short course of oral antibiotics or antibiotic ear drops to be started before or immediately after the procedure, particularly if there is evidence of active infection or a history of recurrent purulent otorrhea.
\end{itemize}

\section*{Contraindications \& Precautions}
While myringotomy and tube insertion is a common and generally safe procedure, certain conditions serve as contraindications or necessitate specific precautions.

\textbf{Absolute Contraindications:}
\begin{itemize}
  \item Active cholesteatoma: The presence of a cholesteatoma, a destructive skin cyst in the middle ear, requires a more extensive and definitive surgical approach (e.g., tympanomastoidectomy) rather than simple tube insertion, which could exacerbate the condition.
  \item Existing tympanic membrane perforation: If a perforation is already present and patent, tube insertion is generally not indicated, as the middle ear is already ventilated. The perforation itself may require surgical repair (tympanoplasty).
  \item Severe, uncorrectable bleeding diathesis: Uncontrolled or severe bleeding disorders that cannot be adequately managed preoperatively pose an unacceptable risk for any surgical procedure, including myringotomy.
\end{itemize}

\textbf{Relative Contraindications \& Precautions:}
\begin{itemize}
  \item Coagulopathy or use of anticoagulants: Patients with underlying coagulopathies or those on anticoagulant/antiplatelet medications require careful evaluation and management. These medications may need to be temporarily discontinued or bridged to minimize the risk of intraoperative or postoperative bleeding.
  \item Immunocompromised state: Patients with compromised immune systems (e.g., HIV, chemotherapy, chronic steroid use) may be at a higher risk for postoperative infection and require enhanced prophylactic measures and close monitoring.
  \item Anatomical abnormalities of the external auditory canal or middle ear: Severe stenosis of the ear canal, exostoses, or significant middle ear malformations can make tube insertion technically challenging or increase the risk of complications.
  \item Uncontrolled systemic illness: Severe, uncontrolled medical conditions such as decompensated cardiac disease, uncontrolled diabetes, or active respiratory infections may increase anesthetic or surgical risks and should ideally be optimized before elective surgery.
  \item Patient or parent non-compliance: Inability or unwillingness to adhere to postoperative care instructions, particularly regarding water precautions and follow-up appointments, can significantly increase the risk of complications and compromise the success of the procedure.
  \item Water precautions: Patients with tympanostomy tubes should generally avoid direct water entry into the ear canal, especially in lakes, rivers, or non-chlorinated pools, unless protective earplugs are used. Showering and bathing are usually safe without earplugs if the head is not submerged.
  \item Flying and diving: While tubes equalize pressure, preventing barotrauma during air travel, deep diving should be avoided as the increased pressure can force water into the middle ear through the tube.
\end{itemize}

\section*{Complications \& Risks}
While generally considered a safe procedure with a low complication rate, myringotomy with tube insertion carries several potential risks, which can be categorized as general surgical and procedure-specific.

\begin{itemize}
  \item \textbf{Anesthesia Risks:}
    \begin{itemize}
      \item Adverse reactions to anesthetic medications, including nausea, vomiting, allergic reactions, or drug interactions.
      \item Respiratory complications such as laryngospasm, bronchospasm, or aspiration, particularly in children.
      \item Cardiovascular events, though extremely rare in healthy pediatric populations undergoing minor procedures.
    \end{itemize}
  \item \textbf{General Surgical Risks:}
    \begin{itemize}
      \item Bleeding: Usually minimal and self-limiting, but can occasionally require intervention.
      \item Infection: Postoperative otitis externa (external ear canal infection) or otitis media, although the tubes are intended to reduce middle ear infections.
      \item Pain: Mild postoperative pain is common but typically well-controlled with over-the-counter analgesics.
    \end{itemize}
  \item \textbf{Procedure-Specific Risks:}
    \begin{itemize}
      \item \textbf{Otorrhea (Ear Drainage):} This is the most common complication, occurring in 10-20\% of patients. It is usually due to infection and often responds effectively to topical antibiotic ear drops.
      \item \textbf{Premature Extrusion:} The tympanostomy tube may fall out earlier than expected (e.g., before 6 months), leading to a recurrence of symptoms and potentially necessitating re-insertion.
      \item \textbf{Tube Obstruction:} The tube can become blocked by blood clots, cerumen, or middle ear secretions, rendering it ineffective in ventilating the middle ear.
      \item \textbf{Persistent Perforation:} After the tube spontaneously extrudes, the myringotomy site may fail to close, leaving a small, persistent hole in the eardrum. This occurs in 1-3\% of cases and may require surgical repair (tympanoplasty) in the future.
      \item \textbf{Tympanosclerosis:} White, chalky patches of scar tissue may form on the tympanic membrane, which are usually asymptomatic but can rarely affect hearing if extensive.
      \item \textbf{Granulation Tissue:} Small polyps or granulation tissue can form around the tube insertion site, sometimes requiring topical treatment or surgical removal.
      \item \textbf{Cholesteatoma:} Although very rare (less than 0.1\%), a skin cyst can form in the middle ear, potentially requiring more extensive surgery. This risk is slightly higher with long-term tubes or in patients with severe retraction.
      \item \textbf{Hearing Changes:} While the primary goal is to improve hearing, rarely, the procedure itself or complications can lead to temporary or, in very rare instances, permanent changes in hearing.
      \item \textbf{Retained Tube:} In some cases, particularly with T-tubes, the tube may not extrude spontaneously and may require surgical removal.
      \item \textbf{Displacement of Ossicles:} Extremely rare, but possible injury to the ossicular chain during insertion, potentially affecting hearing.
    \end{itemize}
\end{itemize}

\section*{Postoperative Recovery Timeline}
Immediately after myringotomy and tube insertion, patients are transferred to the Post-Anesthesia Care Unit (PACU) for close monitoring as they recover from anesthesia. Children may experience initial drowsiness, irritability, or mild nausea. Pain is typically minimal and can be effectively managed with over-the-counter analgesics such as acetaminophen or ibuprofen. Most patients are discharged home within a few hours of the procedure. It is common to observe a small amount of clear or blood-tinged drainage from the ear canal for the first 24-48 hours, which usually resolves spontaneously. Hearing improvement is often immediate, as any accumulated middle ear fluid has been removed and the middle ear is now properly ventilated.

Over the first 1-2 weeks, patients are advised to avoid strenuous activities, heavy lifting, and activities that could introduce water into the ear canal. Specific instructions regarding ear care, including water precautions (e.g., using earplugs for swimming or avoiding head submersion during bathing), are provided. Antibiotic ear drops may be prescribed for a few days to prevent infection, especially if purulent fluid was drained. The tympanostomy tubes typically remain in place for an average of 6 to 18 months, depending on the tube type, before spontaneously extruding into the ear canal. In the vast majority of cases, the myringotomy site on the eardrum heals completely after the tube falls out. Regular follow-up appointments with the otolaryngologist are crucial to monitor the tubes, assess hearing, and address any potential complications.

\section*{Surgical Findings \& Pathology}
During myringotomy and tube insertion, the surgeon meticulously documents several key findings. These include the appearance of the tympanic membrane (e.g., retracted, bulging, inflamed), the presence and character of middle ear fluid (e.g., serous, mucoid, purulent, thick, thin), and the ease of fluid aspiration. The specific quadrant of the myringotomy incision and the type and size of the tympanostomy tube inserted are also recorded. The patency of the tube and successful aeration of the middle ear are confirmed visually. While the procedure itself does not typically involve sending tissue for routine histopathological examination, any aspirated middle ear fluid may be sent for microbiological culture and sensitivity testing, particularly in cases of recurrent purulent otorrhea, immunocompromised patients, or when unusual pathogens are suspected. The pathology report, if fluid is sent, would detail the presence of bacteria, fungi, or other microorganisms and their susceptibility to various antibiotics, guiding targeted antimicrobial therapy.

\section{section*{Expected Postoperative Course}}
A normal and successful postoperative course following myringotomy and tube insertion is characterized by a rapid resolution of the primary symptoms that led to the surgery. Patients typically experience an immediate improvement in hearing, as the middle ear is now aerated and free of fluid. Any preoperative ear pain or pressure usually subsides quickly. The frequency of acute ear infections should significantly decrease, leading to fewer antibiotic prescriptions and an improved quality of life for the patient and family. Mild, clear or bloody ear drainage may occur for a day or two, but persistent or purulent otorrhea is not expected and warrants medical evaluation. The tympanostomy tubes are expected to remain patent and in place for their intended duration, providing continuous middle ear ventilation. Over time, as the Eustachian tube function matures, the tubes will spontaneously extrude, and the tympanic membrane should heal completely, leaving minimal or no residual scarring.

\section{section*{Postoperative Care \& Follow-up}}
Comprehensive postoperative care and regular follow-up are essential to ensure the success of myringotomy and tube insertion and to manage any potential complications.

\begin{itemize}
  \item \textbf{Postoperative Visits:} The initial follow-up visit is typically scheduled 2-4 weeks after surgery to confirm tube placement and patency. Subsequent visits are usually every 3-6 months until the tubes spontaneously extrude.
  \item \textbf{Otoscopic Examination:} At each visit, the otolaryngologist performs a detailed otoscopic examination to verify that the tubes are patent, properly positioned, and that the middle ear remains clear of fluid.
  \item \textbf{Audiometry and Tympanometry:} Repeat audiometry may be performed periodically to monitor hearing improvement, especially if there were concerns about speech development or if hearing loss persists. Tympanometry can confirm middle ear aeration.
  \item \textbf{Management of Otorrhea:} If ear drainage occurs, the doctor will assess for infection and typically prescribe topical antibiotic ear drops. Oral antibiotics may be necessary for persistent or severe infections.
  \item \textbf{Tube Extrusion Monitoring:} The natural extrusion of the tubes is monitored. If a tube remains in place for an unusually long time (e.g., >2 years for short-term tubes), surgical removal may be considered to prevent complications.
  \item \textbf{Repeat Tube Placement:} If recurrent otitis media or chronic effusion returns after the tubes extrude, a second set of tubes may be necessary.
  \item \textbf{Adenoidectomy Consideration:} In some cases, particularly for recurrent AOM in older children, adenoidectomy may be considered concurrently with or subsequent to tube placement, especially if adenoid hypertrophy is contributing to Eustachian tube dysfunction.
\item \textbf{Warning Signs:} Parents or adult patients should be advised to contact the doctor immediately if they notice persistent ear drainage, significant ear pain, fever, sudden hearing changes, or if the tube appears to have fallen out prematurely.
\end{itemize}

\section*{Epidemiology}
Myringotomy with tympanostomy tube insertion is the most common ambulatory surgical procedure performed on children in the United States, with an estimated 600,000 to 700,000 procedures annually. The peak incidence occurs in children between 1 and 3 years of age, reflecting the higher prevalence of otitis media in this age group due to immature Eustachian tube function, frequent upper respiratory infections, and increased exposure in daycare settings. Boys and girls are affected relatively equally. The primary indications are recurrent acute otitis media (AOM) and chronic otitis media with effusion (OME), with OME accounting for a significant proportion of cases, particularly when associated with documented hearing loss and developmental concerns. While the procedure has an extremely low mortality rate, morbidity is primarily related to the complications discussed previously, such as otorrhea (10-20\%), premature tube extrusion, and persistent tympanic membrane perforation (1-3\%). The overall impact on public health is substantial, given the high incidence of otitis media and its potential effects on childhood development, healthcare utilization, and antibiotic consumption.

\section*{Outcomes \& Prognosis}
The outcomes of myringotomy and tympanostomy tube insertion are generally excellent, particularly for the immediate resolution of middle ear effusions and improvement in hearing. For chronic otitis media with effusion, the success rate for immediate fluid resolution and significant hearing improvement exceeds 90\%. For recurrent acute otitis media, tube insertion significantly reduces the frequency of infections, often by 50-70\% during the period the tubes are in place, leading to a substantial decrease in antibiotic use. The prognosis for most children is very good, with the majority experiencing resolution of their middle ear issues as they grow older and their Eustachian tube function matures. Approximately 20-30\% of children may require a second set of tubes, and a smaller percentage may need three or more sets, particularly those with underlying conditions like cleft palate. Long-term studies indicate that while tympanosclerosis (scarring of the eardrum) is a common finding after tube extrusion, it rarely causes significant long-term hearing impairment. Persistent tympanic membrane perforation requiring surgical repair occurs in a small minority of patients (1-3\%). Overall, the procedure has a positive impact on speech and language development, academic performance, and quality of life, while reducing the need for repeated antibiotic courses. Prognostic factors for recurrence include younger age at initial surgery, presence of craniofacial anomalies, and continued exposure to environmental risk factors such as daycare attendance and passive smoke.

\section*{References}
\begin{enumerate}
  \item MedlinePlus. Myringotomy and tubes. U.S. National Library of Medicine; 2024. Available from: \url{https://medlineplus.gov/ency/article/003016.htm}
  \item Kliegman RM, St. Geme JW, Blum NJ, Shah SS, Tasker RC, Wilson KM, editors. Nelson Textbook of Pediatrics. 22nd ed. Philadelphia, PA: Elsevier; 2024.
  \item Rosenfeld RM, Shin JJ, Schwartz SR, et al. Clinical Practice Guideline: Otitis Media with Effusion (Update). Otolaryngology--Head and Neck Surgery. 2016;154(1\_suppl):S1-S41.
  \item Flint PW, Haughey BJ, Lund VJ, et al, editors. Cummings Otolaryngology: Head and Neck Surgery. 8th ed. Philadelphia, PA: Elsevier; 2021.
\end{enumerate}

\clearpage